\documentclass[12 pt, letterpaper]{hmcpset}
\usepackage{hw-preamble}
\usepackage{pdflscape}

\name{Jayden Thadani}
\class{MATH 145 --- Topics in Geometry and Topology}
\assignment{Homework 11}
\duedate{23rd April 2025}

\begin{document}

\begin{problem}[59.]
	Recall that the $n$-orthoplex is defined to be the simplicial complex with $2n$ vertices
	\[
		V = \{ \pm 1, \dots, \pm n \}
	\]
	whose simplices are the subsets that do not contain any pair $\pm k$. It is a model for the sphere $S^{n - 1}$.

	\begin{enumerate}
		\item List the vertices, edges, and triangles of the $3$-orthoplex.

		\item Write down the matrices for the boundary maps $\partial_{0}, \partial_{1}, \partial_{2}$ and determine their ranks.

		\item Obtain the Betti sequence for the $3$-orthoplex.
	\end{enumerate}
\end{problem}

\begin{solution}
	Denote the $3$-orthoplex by $X$.
	\begin{enumerate}
		\item
			\[
				\boxed{
					\begin{split}
						X^{(0)} & = \{ \{ \pm 1 \}, \{ \pm 2 \}, \{ \pm 3 \} \} \\
						X^{(1)} \setminus X^{(0)} & = \{ \{ \pm 1, \pm 2 \}, \{ \pm 2, \pm 3 \}, \{ \pm 3, \pm 1 \} \} \\
						X^{(2)} \setminus X^{(1)} & = \{ \{ \pm 1, \pm 2, \pm 3 \} \}
					\end{split}
				}
			\]

			We can write the $k$-simplices as $X^{(k)} \setminus X^{(k - 1)}$. We get the sets above by considering all subsets of size $k - 1$ that contain no pairs $\pm m$.

		\item
			\[
				\boxed{
					\begin{split}
						\setcounter{MaxMatrixCols}{20}
						\partial_{0} & = \begin{pmatrix}
1 &  0 &  1 &  0 &  0 &  0 &  0 &  0 & -1 & -1 &  0 &  0 \\
 0 &  1 &  0 &  1 &  0 &  0 &  0 &  0 &  0 &  0 & -1 & -1 \\
-1 & -1 &  0 &  0 &  1 &  0 &  1 &  0 &  0 &  0 &  0 &  0 \\
 0 &  0 & -1 & -1 &  0 &  1 &  0 &  1 &  0 &  0 &  0 &  0 \\
 0 &  0 &  0 &  0 & -1 & -1 &  0 &  0 &  1 &  0 &  1 &  0 \\
 0 &  0 &  0 &  0 &  0 &  0 & -1 & -1 &  0 &  1 &  0 &  1 
\end{pmatrix} \\
						\operatorname{rank} \partial_{0} & = 5 \\
						\partial_{1} & = \begin{pmatrix}
1 & 0 & 0 & 1 & 0 & 0 & 0 & 0 \\
0 & 1 & 0 & 0 & 0 & 1 & 0 & 0 \\
0 & 0 & 1 & 0 & 1 & 0 & 0 & 0 \\
0 & 0 & 0 & 0 & 0 & 0 & 1 & 1 \\
1 & 1 & 0 & 0 & 0 & 0 & 0 & 0 \\
0 & 0 & 1 & 0 & 0 & 0 & 1 & 0 \\
0 & 0 & 0 & 1 & 0 & 1 & 0 & 0 \\
0 & 0 & 0 & 0 & 1 & 0 & 0 & 1 \\
1 & 0 & 1 & 0 & 0 & 0 & 0 & 0 \\
0 & 0 & 0 & 1 & 1 & 0 & 0 & 0 \\
0 & 1 & 0 & 0 & 0 & 0 & 1 & 0 \\
0 & 0 & 0 & 0 & 0 & 1 & 0 & 1 
\end{pmatrix} \\
							\operatorname{rank} \partial_{1} & = 7 \\
							\partial_{2} & = 0 \\
							\operatorname{rank} \partial_{2} & = 0.
					\end{split}
				}
			\]

			Let the $0$-chains be given by the following basis
			\[
				C_{0}(X) = \operatorname{span}([1], [-1], [2], [-2], [3], [-2]),
			\]
			the $1$-chains by
			\[
				C_{1}(X) = \operatorname{span}([1, 2], [-1, 2], [1, -2], [-1, -2], [2, 3], \dots, [-2, -3],  [3, 1], \dots, [-3, -1]),
			\]
			and the $2$-chains by
			\[
				C_{2}(X) = \operatorname{span}([1, 2, 3], [-1, 2, 3], [1, -2, 3], [1, 2, -3], [1, -2, -3] [-1, 2, -3], [-1, -2, 3], [-1, -2, -3]).
			\]

			Then we can compute the matrices using the definition of the boundary map and obtain their ranks using computer tools.

			Note that since $X$ is $2$-dimensional, we have $C_{3}(X) = 0$, and thus, $\partial_{2} : C_{3}(X) \to C_{2}(X)$ is the zero map.

		\item
			\[
				\boxed{
					\mathbf{b} = (1, 0, 1, 0, 0, \dots)
				}
			\]

			$X$ clearly has one connected component. Recalling that the kernel of $\partial_{-1}$ is all of $C_{0}(X)$, we can also write this as
			\begin{align*}
				b_{0} & = \dim \ker \partial_{-1} - \dim \operatorname{img} \partial_{0} \\
				      & = \dim C_{0}(X) - \operatorname{rank} \partial_{0} \\
				      & = 6 - 5 \\
				      & = 1.
			\end{align*}

			$X$ has no $1$-holes. This is reflected in the fact that
			\begin{align*}
				b_{1} & = \dim \ker \partial_{0} - \dim \operatorname{img} \partial_{1} \\
				      & = (\dim C_{1}(X) - \operatorname{rank} \partial_{0}) - \operatorname{rank} \partial_{1} \\
				      & = 12 - 5 - 7 \\
				      & = 0.
			\end{align*}

			Finally, $X$ has one $2$-hole, given by
			\begin{align*}
				b_{2} & = \dim \ker \partial_{1} - \dim \operatorname{img} \partial_{2} \\
				      & = (\dim C_{2}(X) - \operatorname{rank} \partial_{1}) - 0 \\
				      & = 8 - 7 \\
				      & = 1.
			\end{align*}
			Note that in the intermediate steps of these computations we use the rank-nullity theorem.
	\end{enumerate}
\end{solution}

\pagebreak

\begin{problem}[60.]
	Recall that $\partial B^{n}$ is defined to be the simplicial complex with $n + 1$ vertices
	\[
		V = \{ 0, 1, \dots, n \}
	\]
	whose simplices are the subsets of cardinality at most $n$. It is a model for the sphere $S^{n - 1}$.

	\begin{enumerate}
		\item List the vertices, edges, triangles, and tetrahedra of $\partial B^{4}$.

		\item Write down the matrices for the boundary maps $\partial_{0}, \partial_{1}, \partial_{2}$ and determine their ranks.

		\item Obtain the Betti sequence for $\partial B^{4}$.
	\end{enumerate}
\end{problem}

\begin{solution}
	Let $X = \partial B^{4}$
	\begin{enumerate}
		\item
			\[
				\boxed{
					\begin{split}
						X^{(0)} & = \{ \{ 0 \}, \{ 1 \}, \{ 2 \}, \{ 3 \}, \{ 4 \} \} \\
						X^{(1)} \setminus X^{(0)} & = \{ \{ 0, 1 \}, \{ 0, 2 \}, \{ 0, 3 \}, \{ 0, 4 \}, \{ 1, 2 \}, \{ 1, 3 \}, \{ 1, 4 \}, \{ 2, 3 \}, \{ 2, 4 \}, \{ 3, 4 \} \} \\
						X^{(2)} \setminus X^{(1)} & = \{ \{ 2, 3, 4 \}, \{ 1, 3, 4 \}, \{ 1, 2, 4 \}, \{ 1, 2, 3 \}, \{ 0, 3, 4 \}, \{ 0, 2, 4 \}, \{ 0, 2, 3 \}, \{ 0, 1, 4 \}, \{ 0, 1, 3 \}, \{ 0, 1, 2 \} \} \\
						X^{(3)} \setminus X^{(2)} & = \{ \{ 1, 2, 3, 4 \}, \{ 0, 2, 3, 4 \}, \{ 0, 1, 3, 4 \}, \{ 0, 1, 2, 4 \}, \{ 0, 1, 2, 3 \} \}
					\end{split}
				}
			\]

			Again, we write the set of all $k$-simplices as the difference of skeletons $X^{(k)} \setminus X^{(k - 1)}$.

		\item
			\[
				\boxed{
					\begin{split}
						\partial_{0} & = \begin{pmatrix}
1 &  1 &  1 &  1 &  0 &  0 &  0 &  0 & 0 &  0 \\
-1 &  0 &  0 &  0 &  1 &  1 &  1 &  0 & 0 &  0 \\
 0 & -1 &  0 &  0 & -1 &  0 &  0 &  1 & 1 &  0 \\
 0 &  0 & -1 &  0 &  0 & -1 &  0 & -1 & 0 &  1 \\
 0 &  0 &  0 & -1 &  0 &  0 & -1 &  0 & -1 & -1 
\end{pmatrix} \\
						\operatorname{rank} \partial_{0} & = 4 \\
						\partial_{1} & = \begin{pmatrix}
 0 &  0 &  0 &  0 &  0 &  0 &  0 &  1 &  1 &  1 \\
 0 &  0 &  0 &  0 &  0 &  1 &  1 &  0 &  0 & -1 \\
 0 &  0 &  0 &  0 &  1 &  0 & -1 &  0 & -1 &  0 \\
 0 &  0 &  0 &  0 & -1 & -1 &  0 & -1 &  0 &  0 \\
 0 &  0 &  1 &  1 &  0 &  0 &  0 &  0 &  0 &  1 \\
 0 &  1 &  0 & -1 &  0 &  0 &  0 &  0 &  1 &  0 \\
 0 & -1 & -1 &  0 &  0 &  0 &  0 &  1 &  0 &  0 \\
 1 &  0 &  0 &  1 &  0 &  0 &  1 &  0 &  0 &  0 \\
-1 &  0 &  1 &  0 &  0 &  1 &  0 &  0 &  0 &  0 \\
 1 &  1 &  0 &  0 &  1 &  0 &  0 &  0 &  0 &  0 \\
\end{pmatrix} \\
						\operatorname{rank} \partial_{1} & = 6 \\
						\partial_{2} & = \begin{pmatrix}
1 &  1 &  0 &  0 &  0 \\
-1 &  0 &  1 &  0 &  0 \\
 1 &  0 &  0 &  1 &  0 \\
-1 &  0 &  0 &  0 &  1 \\
 0 & -1 & -1 &  0 &  0 \\
 0 &  1 &  0 & -1 &  0 \\
 0 & -1 &  0 &  0 & -1 \\
 0 &  0 &  1 &  1 &  0 \\
 0 &  0 & -1 & 0 &  1 \\
 0 &  0 &  0 &  -1 & -1 \\
\end{pmatrix} \\
						\operatorname{rank} \partial_{2} & = 5.
					\end{split}
				}
			\]

			Let us have the following bases
			\begin{align*}
				C_{0}(X) & = \operatorname{span}([1], \dots, [4]) \\
				C_{1}(X) & = \operatorname{span}([ 0, 1 ], [ 0, 2 ], [ 0, 3 ], [ 0, 4 ], [ 1, 2 ], [ 1, 3 ], [ 1, 4 ], [ 2, 3 ], [ 2, 4 ], [ 3, 4 ]) \\
				C_{2}(X) & = \operatorname{span}([ 2, 3, 4 ], [ 1, 3, 4 ], [ 1, 2, 4 ], [ 1, 2, 3 ], [ 0, 3, 4 ], [ 0, 2, 4 ], [ 0, 2, 3 ], [ 0, 1, 4 ], [ 0, 1, 3 ], [ 0, 1, 2 ]). \\
				C_{3}(X) & = \operatorname{span}([1, 2, 3, 4], [0, 2, 3, 4], [0, 1, 3, 4], [0, 1, 2, 4], [0, 1, 2, 3])
			\end{align*}
			Again, this leads to exactly one way to write down the matrices from the formula for the boundary map. We obtain the matrices above. The ranks are obtained using computer software.

			Similarly to the previous problem, we have $\partial_{3} = 0$ since $\partial B^{4}$ is $3$-dimensional. Thus, it has $C_{4} = 0$ and $\partial_{3} : C_{4} \to C_{3}$ is really $\partial_{3} = 0$.

		\item
			\[
				\boxed{
					\mathbf{b} = (1, 0, 0, 1, 0, \dots)
				}
			\]
			$X$ clearly has one connected component. Recalling that the kernel of $\partial_{-1}$ is all of $C_{0}(X)$, we can also write this as
			\begin{align*}
				b_{0} & = \dim \ker \partial_{-1} - \dim \operatorname{img} \partial_{0} \\
				      & = \dim C_{0}(X) - \operatorname{rank} \partial_{0} \\
				      & = 5 - 4 \\
				      & = 1.
			\end{align*}

			$X$ has no $1$-holes. This is reflected in the fact that
			\begin{align*}
				b_{1} & = \dim \ker \partial_{0} - \dim \operatorname{img} \partial_{1} \\
				      & = (\dim C_{1}(X) - \operatorname{rank} \partial_{0}) - \operatorname{rank} \partial_{1} \\
				      & = 10 - 4 - 6 \\
				      & = 0.
			\end{align*}

			$X$ also has no $2$-holes. This is reflected in the fact that
			\begin{align*}
				b_{3} & = \dim \ker \partial_{1} - \dim \operatorname{img} \partial_{2} \\
				      & = (\dim C_{2}(X) - \operatorname{rank} \partial_{1}) - \operatorname{rank} \partial_{2} \\
				      & = 10 - 6 - 4 \\
				      & = 0.
			\end{align*}

			Finally, $X$ has one $3$-hole, given by
			\begin{align*}
				b_{2} & = \dim \ker \partial_{2} - \dim \operatorname{img} \partial_{3} \\
				      & = (\dim C_{3}(X) - \operatorname{rank} \partial_{2}) - 0 \\
				      & = 5 - 4 - 0 \\
				      & = 1.
			\end{align*}
			In the intermediate steps of these computations we use the rank-nullity theorem.
	\end{enumerate}
\end{solution}

\pagebreak

\begin{problem}[61. (the Morse inequalities)]
	Prove the following inequality for a finite simplicial complex:
	\[
		b_{k} - b_{k - 1} + \dots + (-1)^{k} b_{0} \le c_{k} - c_{k - 1} + \dots + (-1)^{k} c_{0}.
	\]
\end{problem}

\begin{solution}
	Recall that for the simplicial chain complex
	\[
	\begin{tikzcd}
		0 & {C}_{0}(X, \mathbb{F}) \ar[l] & {C}_{1}(X, \mathbb{F}) \ar[l, "\partial_{0}"] & {C}_{2}(X, \mathbb{F}) \ar[l, "\partial_{1}"] & \dots \ar[l, "\partial_{2}"]
	\end{tikzcd}
	\]
	we have
	\[
		b_{i} = \dim (\ker \partial_{i - 1} / \operatorname{img} \partial_{i}) = \operatorname{null} \partial_{i - 1} - \operatorname{rank} \partial_{i}
	\]
	by definition, and
	\[
		c_{i} = \operatorname{rank} \partial_{i - 1} + \operatorname{null} \partial_{i - 1}
	\]
	by the rank-nullity theorem.

	Note that
	\begin{align*}
		b_{k} - b_{k - 1} & = \operatorname{null} \partial_{k - 1} - \operatorname{rank} \partial_{k} + \operatorname{null} \partial_{k - 2}
	\end{align*}
	and
	\begin{align*}
		b_{i} - \operatorname{null} \partial_{i} & = \operatorname{null} \partial_{i - 1} - \operatorname{rank} \partial_{i} - \operatorname{null} \partial_{i} \\
							 & = \operatorname{null} \partial_{i - 1} - c_{i}.
	\end{align*}
	We can use this to turn the sum into something like a telescoping sum. We have
	\begin{align*}
		\sum_{i = 0}^{k} (-1)^{k - i} b_{i} & = - \operatorname{rank} \partial_{k} + c_{k}  - \operatorname{null} \partial_{k - 2} + \sum_{i = 0}^{k - 2} (-1)^{k - i} b_{i} \\
						& = - \operatorname{rank} \partial_{k} + c_{k} + b_{k - 2} - \operatorname{null} \partial_{k - 2} + \sum_{i = 0}^{k - 3} (-1)^{k - i} b_{i} \\
						& = -\operatorname{rank} \partial_{k} + c_{k} - c_{k - 1} + \operatorname{null} \partial_{k - 3} + \sum_{i = 0}^{k - 3} \\
						& = \vdots \\
						& = -\operatorname{rank} \partial_{k} + \sum_{i = 0}^{k} (-1)^{k - i} c_{k}.
	\end{align*}

	That is,
	\[
		\sum_{i = 0}^{k} (-1)^{k - i} b_{i} = - \operatorname{rank} \partial_{k} + \sum_{i = 0}^{k} (-1)^{k - i} c_{i} \le \sum_{i = 0}^{k} (-1)^{k - i} c_{i}.
	\]
	If $X$ is $k$-dimensional, then all $n > k$ have $C_{n}(X, \mathbb{F}) = 0$. In particular $C_{k + 1}(X, \mathbb{F}) = 0$ and thus, $\partial_{k} : C_{k + 1}(X, \mathbb{F}) \to C_{k}(X, \mathbb{F})$ has $\operatorname{rank} \partial_{k} = 0$. This is exactly Euler's formula.
\end{solution}

\pagebreak

\begin{problem}[62.]
	The torus is a topological space which can be constructed from a closed square by gluing each side to its opposite side. Here is a representation of a torus as a simplicial complex with $9$ vertices, $27$ edges, and $18$ triangles.
	\begin{figure}[H]
		\centering
		\includegraphics[width = 0.3 \textwidth]{figures/P62 torus}
	\end{figure}

	By calculating the ranks of $\partial_{0}, \partial_{1}$, determine the Betti sequence for the torus.
\end{problem}

\begin{solution}
	\[
		\boxed{
			\mathbf{b} = (1, 2, 1, 0, 0, \dots)
		}
	\]

	We denote the torus by $\mathbb{T}$.

	We used the bases
	\begin{align*}
		C_{0}(X) & = \operatorname{span}([1], [2], \dots, [9]) \\
		C_{1}(X) & = \operatorname{span}([1, 4], [4, 7], [7, 1], [2, 1], [5, 4], [8, 7], [2, 4], [5, 7], [8, 1], [2, 5], [5, 8], [8, 2], [3, 2] \\
			 & [6, 5], [9, 8], [3, 5], [6, 8], [9, 2], [3, 6], [6, 9], [9, 3], [1, 3], [4, 6], [7, 9], [1, 6], [4, 9], [7, 3]) \\
		C_{2}(X) & = \operatorname{span}([2, 1, 4], [5, 4, 7], [8, 7, 1], [2, 5, 4], [5, 8, 7], [8, 2, 1], [3, 2, 5], [6, 5, 8], [9, 8, 2], \\
			 & [3, 6, 5], [6, 9, 8], [9, 3, 2], [3, 6, 1], [4, 6, 9], [7, 9, 3], [1, 4, 6], [4, 7, 9], [7, 1, 3]).
	\end{align*}

	Computing the ranks of the matrices obtained by applying the boundary maps to these basis vectors, we get $\operatorname{rank} \partial_{0} = 8$ and $\operatorname{rank} \partial_{1} = 17$. Also, again because $\mathbb{T}$ is $2$-dimensional, we have $\partial_{2} = 0$ and all higher (than $2$) homology vanishes.

	$\mathbb{T}$ clearly has one connected component. Recalling that the kernel of $\partial_{-1}$ is all of $C_{0}(X)$, we can also write this as
			\begin{align*}
				b_{0} & = \dim \ker \partial_{-1} - \dim \operatorname{img} \partial_{0} \\
				      & = \dim C_{0}(X) - \operatorname{rank} \partial_{0} \\
				      & = 9 - 8 \\
				      & = 1.
			\end{align*}

			$\mathbb{T}$ has no $1$-holes. This is reflected in the fact that
			\begin{align*}
				b_{1} & = \dim \ker \partial_{0} - \dim \operatorname{img} \partial_{1} \\
				      & = (\dim C_{1}(X) - \operatorname{rank} \partial_{0}) - \operatorname{rank} \partial_{1} \\
				      & = 27 - 8 - 17 \\
				      & = 2.
			\end{align*}

			Finally, $\mathbb{T}$ has one $2$-hole, given by
			\begin{align*}
				b_{2} & = \dim \ker \partial_{1} - \dim \operatorname{img} \partial_{2} \\
				      & = (\dim C_{2}(X) - \operatorname{rank} \partial_{1}) - 0 \\
				      & = 18 - 17 \\
				      & = 1.
			\end{align*}
			Again in the intermediate steps of these computations we use the rank-nullity theorem.

\end{solution}

\pagebreak

\begin{problem}[63.]
	Let $X = \{ 2, 3, 4, \dots, 16 \}$. Define a simplicial complex $S_{16}$ as follows: a subset $\sigma \subseteq X$ belongs to $S_{16}$ if and only if the elements of $\sigma$ share a common factor greater than $1$.
	\begin{enumerate}
		\item Explain why $S_{16}$ is a simplicial complex.

		\item How many vertices and edges in $S_{16}$?

		\item What is the largest-dimensional simplex in $S_{16}$?

		\item How many connected components does $S_{16}$ have?

		\item Find three edges $\{ a, b \}, \{ b, c \}, \{ c, a \}$ in $S_{16}$ such that $\{ a, b, c \} \notin S_{16}$.

		\item What are the \emph{maximal} simplices of $S_{16}$.

		\item Sketch $S_{16}$ and write down your best guess for its Betti sequence.

		\item The maximal simplices form a cover of the set $X$. Determine and sketch the nerve of this cover.

		\item Calculate the Betti sequence of the nerve. Compare it with your guess in (g).
	\end{enumerate}
\end{problem}

\begin{solution}
	\begin{enumerate}
		\item Suppose $\sigma \in S_{16}$ with all elements sharing a common factor $d > 1$ --- we have $d \mid x \in \sigma$ for every $x \in \sigma$. Suppose $\tau \subseteq \sigma$. Then $d \mid x$ for every $x \in \tau$. Thus, all elements of $\tau$ share a common factor greater than $1$.

			This suffices to show $S_{16}$ is a simplicial complex.

		\item 
			\[
				\boxed{
					\begin{split}
						S_{16}^{(0)} & = \{ \{ 2 \}, \{ 3 \}, \dots, \{ 16 \} \} \\
						S_{16}^{(1)} \setminus S_{16}^{(0)} & = \binom{X}{2} \setminus \\
										    & \left (\{ \{ n, p \} \mid p \text{ prime} \} \cup \{ \{ 4, 9 \}, \{ 4, 15 \}, \{ 8, 9 \}, \{ 8, 15 \}, \{ 9, 10 \}, \{ 9, 14 \}, \{ 9, 16 \}, \{ 14, 15 \}, \{ 15, 16 \} \} \right ) . \\
						S_{16}^{(2)} \setminus S_{16}^{(1)} & = \{ \{ m_{1}, m_{2}, m_{3} \} \mid m_{1}, m_{2}, m_{3} \in 2 \mathbb{Z} \} \cup \{ \{ m_{1}, m_{2}, m_{3} \} \mid m_{1}, m_{2}, m_{3} \in 3 \mathbb{Z} \} \cup \{ \{ 5, 10, 15 \} \}
					\end{split}
				}
			\]

			All of the singleton sets are vertices --- they are all greater than one, and are thus, divisible by themselves, a number greater than one.

			For the edges, we can ignore all pairs where one number is prime as well as the pairs of co-prime numbers listed above.

			For triangles we break the triples down by the common prime dividing them. There are $\binom{8}{3}$ divided by $2$, $\binom{5}{3}$ divided by $3$, and $1$ divided by $5$.

		\item
			\[
				\boxed{
					\begin{split}
						\sigma & = \{ 2, 4, 6, \dots, 16 \} \\
						\dim \sigma & = 7.
					\end{split}
				}
			\]

			The largest-dimensional simplex of $S_{16}$ consists of the $8$ even numbers less than $16$.

		\item
			\[
				\boxed{
					b_{0}(X) = 3.
				}
			\]

			Every number that is not a prime greater than $8$ is divisible by a number less than $8$. That number divides an even number. Thus, almost all numbers are connected to the largest simplex of all even numbers.

			The primes greater than $8$, namely $11$ and $13$ share common factors with none of the other numbers, and thus, form two more connected components.

		\item 
			\[
				\boxed{
					\{ 6, 10, 15 \} \notin S_{16}, \text{ but } \{ 6, 10 \}, \{ 10, 15 \}, \{ 15, 6 \} \in S_{16}.
				}
			\]

		\item
			\[
				\boxed{
					\mathcal{U} = \{ U_{p} = \{ pk \mid p \text{ prime}, k \in \mathbb{N} \} \}
				}
			\]
			
			Every simplex $\sigma$ of $S_{16}$ is a set of numbers with a common divisor $d$ greater than $1$. Then the smallest prime $p$ dividing $d$ divides all the numbers in the simplex. Thus, the simplex corresponding to all multiples of $p$, denoted $U_{p}$, contains $\sigma$. 

			These $U_{p}$ are maximal since they include the prime $p$ which is only divisible by itself, and is thus, contained in no larger simplex.

		\item
			\[
				\boxed{
					\mathbf{b} = (3, 1, 0, 0, \dots)
				}
			\]

			\begin{figure}[H]
				\centering
				\includegraphics[width = 0.5 \textwidth]{figures/P63g guess.pdf}
				\caption{$S_{16}$ drawn with $7$ simplices highlighted in red, $5$-simplices highlighted in yellow, $2$-simplices highlighted in black, and edges in black, only drawn outside of simplices}
			\end{figure}

			Clearly this has $3$ connected components.

			Notice that after contracting all the large $k$-simplices onto the drawn edges, we are left with a $1$-cycle $[6, 10] + [10, 15] + [15, 6]$ that is not the boundary of any $2$-chain. Thus, we have exactly one hole.
		\item
			\[
				\boxed{
					\operatorname{nerve} \mathcal{U} = \{ \{ U_{2} \}, \{ U_{3} \}, \{ U_{5} \}, \{ U_{7} \}, \{ U_{11} \}, \{ U_{13} \}, \{ U_{2}, U_{3} \}, \{ U_{2}, U_{5} \}, \{ U_{2}, U_{7} \}  \{ U_{3}, U_{5} \} \}.
				}
			\]
			
			\begin{figure}[H]
				\centering
				\includegraphics[width = 0.5 \textwidth]{figures/P63h nerve.pdf}
				\caption{$\operatorname{nerve} \mathcal{U}$}
			\end{figure}

			Since the condition that two maximal simplicial complexes intersect is equivalent to the condition that their least common multiple is at most $16$, we can write the nerve as
			\[
				\operatorname{nerve} \mathcal{U} = \{ \{ U_{p_{1}}, \dots, U_{p_{n}} \} \mid p_{1} \cdots p_{n} \le 16 \}.
			\]
			(The least common multiple of distinct primes is just their product). There are no triples of primes such that this is true (the three smallest primes have product $2 \times 3 \times 5 = 30 > 16$). The pairs of primes for which this is true are $2 \times 3 = 6$, $2 \times 5 = 10$, $2 \times 7 = 14$ and $3 \times 5 = 15$. Thus,
			\[
				\operatorname{nerve} \mathcal{U} = \{ \{ U_{2} \}, \{ U_{3} \}, \{ U_{5} \}, \{ U_{7} \}, \{ U_{11} \}, \{ U_{13} \}, \{ U_{2}, U_{3} \}, \{ U_{2}, U_{5} \}, \{ U_{2}, U_{7} \}  \{ U_{3}, U_{5} \} \}.
			\]

		\item
			\[
				\boxed{
					\mathbf{b} = (3, 1, 0, 0, \dots)
				}
			\]

			We've already shown that $b_{0} = 3$. In the sketch of the nerve there is only one $1$-cycle, and it's not a $1$-boundary. This is the unique non-trivial dimension of the homology vector space. Thus, $b_{1} = 1$.
			The nerve is $1$-dimensional, so has trivial higher homology.

			The Betti number of the nerve is exactly the same as the estimated Betti number of $S_{16}$.
	\end{enumerate}
\end{solution}

\end{document}
